\section*{Preface}

These notes cover the Linear Algebra course taught during Aug-Nov 2026 at CMI by Prof. Upendra.
Throughout the course the emphasis will be on linear algebra over real and complex fields.

\subsection{Timings}
$3:30-4:45$pm Wednesday, Thursday, Friday. Shared tutorials with Analysis every Monday and
Friday at $2$pm.

\subsection{Textbook}
\emph{Algebra} by Michael Artin, second edition.
Other reference material: \emph{Linear Algebra} by Kenneth Hoffman and Rau Kunze,
\emph{Linear Algebra Done Right} by Sheldon Axler.

\subsection{Syllabus}

Topics covered will be --- Systems of linear equations, row reduction, vector spaces over any field, bases and
dimension, change of basis, linear transformations, dimension formula, matrix of a linear transformation
and change of bases, linear operators, similarity, determinants and invertibility, eigenvalues and eigenspaces,
characteristic polynomial, triangular and diagonal forms, Euclidean/Hermitian spaces with standard inner
product, orthogonal projection, orthonormal bases and Gram-Schmidt procedure, spectral theorem for Her-
mitian/symmetric operators. The above material corresponds to the following from Artin’s Algebra: chapters
$1, 3, 4.1\mbox{--}4.6$, parts of chapter $8$.

\subsection{Grading}

Homework will \emph{not} be graded for the most part. He argues that this is due to LLMs, and this seems reasonable to me.
Loosely, the grading will be based on two-three quizes, two tests, and a final exam. I don't remember
the exact distribution but it was something like $5-10\%$ for homework, $10-20\%$ for the quizzes, $30\%$ for
the tests and $30-40\%$ for the final exam. Attendance will not be part of the grade, but is strongly
encouraged.
